\documentclass{article}
\usepackage{preamble}

\begin{document}
\header{Filemon Mateus}{CS 6140}{Distances \& LSH}{Homework \# 3}{\today}
\begin{enumerate}[leftmargin={*}, font={\bf}, label={\arabic*.}, ref={\arabic*}]
  \item \label{qst:1}
    \begin{enumerate}[label={(\Alph*)}, ref={\theenumi(\Alph*)}]
      \item \label{qst:1A}
        From Section $5.2$ in M4D, given a budget of $t$ hash functions and a threshold
        $\tau$, the value of $r$ is approximately given by:
        \[
          r \approx -\log_{\tau}{(t)} = -\log_{.75}{(200)} \approx 19
        \]
        Furthermore, the number of hash functions $t$ is split into $r$ and $b$; therefore,
        $b$---the number of rows within each band---is the ratio between $t$ and $r$, which
        corresponds to:
        \[
          b = t/r = 200/19 \approx 11
        \]
        We can verify that these values yield a good separation around $\tau$ by plugging
        them into $f(s) = 1 - (1-s^b)^r$ and check for an $S$-shaped curve with steepest
        slope/rise at $\tau=0.75$.
        \begin{figure}[H]
          \centering
          \begin{tikzpicture}
            \begin{axis}[
              xlabel={$\text{JS}(D_1, D_2)$},
              ylabel={Probability of checking $d(D_1, D_2)$},
              ymin={-0.1},
              ymax={+1.1},
            ]
              \addplot[
                thick,
                color=red
              ] table[col sep=comma, x index=0, y index=1] {dev-dir/week3/data/lsh-values.csv};
              \addlegendentry{$f(s) = 1 - (1 - s^b)^r$};
              \addplot[
                thin,
                dashed
              ] coordinates { (0.75, -0.1) (0.75, 1.1) };
              \addlegendentry{$\tau$-value = $0.75$};
            \end{axis}
          \end{tikzpicture}
          \caption{
            Probability that the distance $d(D_1, D_2)$ is checked in a LSH scheme with \\
            $b=11$ bands with $r=19$ hashes each, as function of $s = \text{JS}(D_1, D_2)$.
            See \autoref{lst:lsh-script}.
          }
          \label{fig:lsh-curve}
        \end{figure}

      \item \label{qst:1B}
        The six numbers outlining the probability for each pair of the four objects being
        estimated as similar can be retrieved as follows. First, we note that
        \begin{gather*}
          \text{JS}(A,B) = 0.77 \quad \text{JS}(A,C) = 0.25 \quad \text{JS}(A,D) = 0.33 \\
          \text{JS}(B,C) = 0.20 \quad \text{JS}(B,D) = 0.55 \quad \text{JS}(C,D) = 0.91
        \end{gather*}
        and each one of these values correspond to $s$. Furthermore, for a threshold $\tau=
        0.75$, we want $f(s)$ parameterized with $r=19$ and $b=11$. Hence plugging each $s$
        into $f(s)$, we obtain the following results:
        \begin{table}[H]
          \centering
          \begin{tabular}{|cccc|}
            \hline
            & {\tt B} & {\tt C} & {\tt D} \\ \hline \hline
            {\tt A} & 0.66823 & 0.00000 & 0.00010 \\ \hline
            {\tt B} & --- & 0.00000 & 0.02614 \\ \hline
            {\tt C} & --- & --- & 0.99975 \\ \hline
          \end{tabular}
          \caption{
            Probability that each pair of documents will have similarity $> \tau$.
          }
          \label{tbl:build-ngram-table}
        \end{table}
    \end{enumerate}

  \item \label{qst:2}
    \begin{enumerate}[label={(\Alph*)}, ref={\theenumi(\Alph*)}]
      \item \label{qst:2A}
        From Section $4.6.4$ in M4D, we construct a single random unit vector $\bm{x}
        \in \mathbb{R}^{10}$ (chosen uniformily over the space of all unit vectors) as
        follows:
        \begin{enumerate*}[label={(\roman*)}, itemjoin={,~}, itemjoin*={,~and~}]
          \item
            sample $u_1, u_2 \overset{\mathrm{i.i.d.}}{\sim} \mathcal{U}(0,1)$
          \item
            define $1$-dimensional gaussian random variable $\mathcal{Y} = \sqrt{-2\ln(u_1)}\cos(2\pi u_2)$
          \item
            sample each $x_1, \ldots, x_{10} \overset{\mathrm{i.i.d.}}{\sim} \mathcal{Y}$
          \item
            scale $\bm{x}$ by the inverse of its norm to obtain the desired random unit
            vector $\bar{\bm{x}} = \bm{x} / \|\bm{x}\|_2$.
        \end{enumerate*}

      \item \label{qst:2B}
        The pairwise dot product CDF is illustrated below.
        \begin{figure}[H]
          \centering
          \begin{tikzpicture}
            \begin{axis}[
              xlabel={Pairwise Dot Product},
              ylabel={Likelihood of Occurrence},
            ]
              \addplot [
                thick,
                color=red,
                hist={data=x, bins=1000, handler/.style={const plot}, cumulative, density}
              ] file {dev-dir/week3/data/dot-products.csv};
              \addlegendentry{CDF};
            \end{axis}
          \end{tikzpicture}
          \caption{
            Dot product for each pair of random unit vectors in $\mathbb{R}^{100}$.
            See \autoref{lst:gen-unit-vec}.
          }
          \label{fig:dot-products-cdf-plot}
        \end{figure}
    \end{enumerate}

  \item \label{qst:3}
    \begin{enumerate}[label={(\Alph*)}, ref={\theenumi(\Alph*)}]
      \item \label{qst:3A}
        The pairwise angular similarity CDF is illustrated below. The number of similarities
        that exceed $\tau=0.75$ is {\bf 107004}.
        \begin{figure}[H]
          \centering
          \begin{tikzpicture}
            \begin{axis}[
              xlabel={Pairwise Angular Similarity},
              ylabel={Likelihood of Occurrence},
            ]
              \addplot [
                thick,
                color=red,
                hist={data=x, bins=1000, handler/.style={const plot}, cumulative, density}
              ] file {dev-dir/week3/data/angular-sims-r-dataset.csv};
              \addlegendentry{CDF};
            \end{axis}
          \end{tikzpicture}
          \caption{
            Angular similarities of pairwise unit vectors from {\tt R.txt} dataset. See
            \autoref{lst:angular-sim}.
          }
          \label{fig:angular-sim-r-dataset-cdf-plot}
        \end{figure}

      \item \label{qst:3B}
        The pairwise angular similarity CDF is illustrated below. The number of similarities
        that exceed $\tau=0.75$ is {\bf 0}.
        \begin{figure}[H]
          \centering
          \begin{tikzpicture}
            \begin{axis}[
              xlabel={Pairwise Angular Similarity},
              ylabel={Likelihood of Occurrence},
            ]
              \addplot [
                thick,
                color=red,
                hist={data=x, bins=1000, handler/.style={const plot}, cumulative, density}
              ] file {dev-dir/week3/data/angular-sims.csv};
              \addlegendentry{CDF};
            \end{axis}
          \end{tikzpicture}
          \caption{
            Angular similarities for each pair of random unit vectors in $\mathbb{R}^{100}$.
            See \autoref{lst:gen-unit-vec}.
          }
          \label{fig:angular-sim-cdf-plot}
        \end{figure}
    \end{enumerate}
\end{enumerate}

\newpage

\section*{Appendix}

\lstinputlisting
  [caption={\sf lsh-script.py}, label={lst:lsh-script}]
  {dev-dir/week3/lsh-script.py}
  
\newpage

\lstinputlisting
  [caption={\sf gen-unit-vec.py}, label={lst:gen-unit-vec}]
  {dev-dir/week3/gen-unit-vec.py}
  
\newpage

\lstinputlisting
  [caption={\sf angular-sim.py}, label={lst:angular-sim}]
  {dev-dir/week3/angular-sim.py}
\end{document}
